\section{Introdução}
\label{sec:introducao}

A gestão de edificações públicas universitárias enfrenta uma tensão decisória recorrente: intervenções visíveis ou urgentes competem por recursos com ações preventivas, ambientais, sociais e institucionais cujos benefícios se distribuem ao longo do ciclo de vida. Estudos de caso em instituições universitárias sugerem que a ausência de critérios transparentes favorece práticas reativas e dificulta justificar tecnicamente a prioridade atribuída a cada ativo \parencite{lateef_manutencaouniversidade_2010,barbosa_facilitiesinstituicaopublica_2020}. A manutenção predial condiciona desempenho, segurança, durabilidade e continuidade dos serviços, e a NBR 5674 estabelece requisitos de planejamento, controle e avaliação do sistema de manutenção \parencite{abnt_nbr5674_2012}. No âmbito público brasileiro, essa responsabilidade também se vincula a deveres legais e administrativos do gestor \parencite{araujoneto_manutencaolegislacao_2015}.

Em instituições com infraestrutura envelhecida, restrições orçamentárias e demanda acumulada, ordenar intervenções somente por custo imediato ou urgência pode omitir efeitos sobre desempenho futuro e continuidade operacional. A análise de custos em universidade federal demonstra a utilidade de modelos de previsão para o planejamento orçamentário \parencite{martins_custosmanutencaoses_2024}, enquanto registros de ordens de serviço permitem reconhecer padrões de demanda e apoiar o planejamento de intervenções \parencite{morais_eficienciamanutencao_2023}. Tomados em conjunto, esses estudos sustentam a necessidade de combinar custos e ordens de serviço com condição física, risco, energia, impacto ambiental, experiência dos usuários, conformidade e capacidade institucional \parencite{maia_gestaoinformacaomanutencao_2016,nageli_ciclovida_2019,alfalah_frameworkeducacional_2022}.

A gestão sustentável de facilidades (\textit{facility management}) articula desempenho ambiental, econômico e social com planejamento, operação e gestão dos ativos construídos \parencite{nielsen_sustentabilidadefm_2016,opoku_futurofm_2022,okoro_sustainablefm_2023}. Os Objetivos de Desenvolvimento Sustentável oferecem metas públicas para infraestrutura, cidades e uso responsável de recursos, enquanto a abordagem ambiental, social e de governança (ESG, do inglês \textit{environmental, social and governance}) organiza indicadores aplicáveis aos ativos e serviços prediais \parencite{okoro_sustainablefm_2023,ramantswana_esgfm_2026}. Métodos de apoio multicritério podem explicitar compensações entre esses objetivos, desde que critérios, indicadores, pesos, regras não compensatórias e incertezas sejam documentados \parencite{alfalah_frameworkeducacional_2022,naji_campusdelphi_2023,masmoudi_ahptopsis_2026}.

Apesar dessa convergência potencial, estudos sobre digitalização e métodos decisórios ainda documentam separadamente integração de dados, previsão operacional, critérios de sustentabilidade e escolha multicritério, sem demonstrar de modo recorrente a cadeia completa até a decisão institucional \parencite{deng_bimdigitaltwins_2021,das_iotai_2025,alfalah_frameworkeducacional_2022,masmoudi_ahptopsis_2026}. A lacuna abordada neste artigo é a conexão rastreável entre evidência científica, variável observável e regra pública de decisão, sem confundir frequência bibliográfica com importância normativa.

O objetivo geral é analisar a estrutura bibliométrica e temática da literatura e, com base nessa evidência, especificar indicadores e um fluxo de validação para futura parametrização multicritério da priorização de intervenções em edificações públicas universitárias. A pergunta central (RQ0) é: como os padrões bibliométricos e a síntese temática da literatura podem fundamentar critérios, indicadores e etapas de validação para uma futura aplicação multicritério nesse contexto?

A análise foi orientada originalmente por cinco questões específicas: quais dimensões de sustentabilidade são identificadas (RQ1); quais critérios de priorização são identificados (RQ2); quais métodos de apoio à decisão são identificados (RQ3); em quais contextos de edificação essas abordagens aparecem (RQ4); e quais lacunas são registradas para instituições públicas de ensino superior (RQ5). Como verificação complementar, formulou-se a RQ6: em que medida uma busca de sensibilidade dedicada à inteligência artificial (IA) e ao aprendizado de máquina altera a identificação de métodos informacionais e preditivos aplicados à manutenção e à gestão de edificações? A RQ6 permanece uma análise secundária de cobertura e não redefine retrospectivamente a busca principal nem a contribuição central do artigo.

A contribuição central é uma especificação operacional candidata, composta por matriz de critérios e indicadores e por um fluxo de validação institucional. O mapeamento bibliométrico, a síntese temática e a busca de sensibilidade funcionam como camadas de evidência que justificam essa especificação; não são apresentados como produtos independentes de igual hierarquia nem como validação de um modelo decisório pronto.
